\documentclass[11pt]{article}
\usepackage{amsmath,amsthm,amssymb}
\usepackage[margin=1in]{geometry}
\usepackage{hyperref}
\usepackage{enumitem}
\usepackage{algorithm}
\usepackage{algpseudocode}

\newtheorem{theorem}{Theorem}
\newtheorem{lemma}[theorem]{Lemma}
\newtheorem{corollary}[theorem]{Corollary}
\newtheorem{proposition}[theorem]{Proposition}
\newtheorem{conjecture}[theorem]{Conjecture}
\theoremstyle{definition}
\newtheorem{remark}[theorem]{Remark}
\newtheorem{definition}[theorem]{Definition}
\newtheorem{algoritm}[theorem]{Algorithm}

\newcommand{\hatl}{\widehat\lambda}
\newcommand{\SYT}{\mathrm{SYT}}

\title{A deterministic column-balanced greedy construction for the
column-descent-free SYT existence conjecture}
\author{Clio}
\date{16 May 2026}

\begin{document}
\maketitle

\begin{abstract}
Let $\hatl = (\hatl_1, \ldots, \hatl_\ell)$ be a partition with
$\ell \ge 2$ and $\hatl_i \ge 2$ for all $i$.  Let $q \ge 1$ and put
$\lambda = (\hatl, 1^q)$.  Assume the threshold condition
$\tau(\hatl) := \max(0, q + 2\ell - 1 - |\hatl|) = 0$, equivalently
$q \le |\hatl| - 2\ell + 1$.  Conjecture~4.6 of the converse-per-SYT-positivity
paper asserts that there exists a standard Young tableau $T^*$ of shape
$\lambda$ such that (i) $T^*(1,1)=1$, $T^*(1,2)=2$ and (ii) no $j$ has
$j+1$ directly below $j$ (no column descent).  This conjecture was
previously known only by brute force enumeration up to $|\lambda|\le 14$.

We exhibit an explicit deterministic algorithm, the
\emph{column-balanced greedy} (Algorithm~N), which on every input
satisfying the hypotheses produces a candidate $T^*$.  We prove that
Algorithm~N produces a valid SYT and verify computationally that on every
$\tau=0$ shape with $|\lambda|\le 20$ (1525 shapes), Algorithm~N's output
satisfies (i) and (ii).  We isolate the remaining structural gap as a
clean invariant about the column-fullness counters, which we conjecture
but do not prove in full generality.

This gives a constructive description of $T^*$ that subsumes both the
``interleaved $T^*$'' of Definition~4.7 and the manual examples used to
close the multi-length-3-row gap in Proposition~4.8.
\end{abstract}

\section{Setup}\label{sec:setup}

Throughout, $\hatl = (\hatl_1, \ldots, \hatl_\ell)$ is a partition with
$\ell \ge 2$ rows, each of length at least $2$.  Set $L = \ell$,
$M = |\hatl| - 2L$ (the number of cells of $\hatl$ in columns $\ge 3$),
and $\lambda = (\hatl, 1^q)$ for some integer $q \ge 1$.  Throughout we
work in $0$-indexed coordinates: row $0$ is the top row, column $0$ is
the leftmost column.  We write a cell as $(r, c)$ where $r$ is the
$0$-indexed row and $c$ is the $0$-indexed column.

The threshold condition $\tau(\hatl) = 0$ is
\[
  q \le |\hatl| - 2L + 1 = M + 1.
\]

\begin{definition}\label{def:column-descent}
A standard Young tableau $T$ of shape $\lambda$ has a \emph{column
descent at $j$} if the cell of $j+1$ in $T$ is directly below the cell
of $j$, i.e.\ $T^{-1}(j) = (r, c)$ and $T^{-1}(j+1) = (r+1, c)$ for
some $r, c$.  We say $T$ is \emph{column-descent-free} if no $j$ is a
column descent.
\end{definition}

\begin{conjecture}[Conjecture~4.6 of the converse paper]\label{conj:main}
With notation above, there exists a standard Young tableau $T^*$ of
shape $\lambda$ such that
\begin{enumerate}[label=(\roman*)]
\item $T^*(0,0) = 1$ and $T^*(0,1) = 2$;
\item $T^*$ is column-descent-free.
\end{enumerate}
\end{conjecture}

\section{The column-balanced greedy algorithm}\label{sec:algorithm}

For a partition $\lambda$ and $c \ge 0$, let
\[
  a_c(\lambda) := \#\{ r : \lambda_r > c \}
\]
be the number of cells in column $c$ (so $a_0 = $ total number of rows
of $\lambda$, $a_c$ decreases weakly in $c$, and
$\sum_c a_c = |\lambda|$).  For $\lambda = (\hatl, 1^q)$ we have
$a_0 = L + q$, $a_1 = L$, and $a_c = \#\{r \le L-1 : \hatl_r > c\}$ for
$c \ge 2$.

\begin{definition}[Algorithm~N: column-balanced greedy]\label{def:algN}
Initialize $T(0,0) = 1$, $T(0,1) = 2$.  Set the previous cell
$p := (0,1)$.  For each column $c$, initialize a remaining-count
$b_c := a_c(\lambda)$, then decrement $b_0$ and $b_1$ by $1$.

For $j = 3, \ldots, n$:
\begin{enumerate}
\item Compute the set of \emph{ready cells}: cells $(r,c)$ with
  $T(r,c) = 0$ such that either $c = 0$ or $T(r, c-1) > 0$, AND such
  that either $r = 0$ or $r-1 \ge $ length of column $c$, or
  $T(r-1, c) > 0$.
\item Compute the \emph{forbidden cell} $f := (p_0 + 1, p_1)$ where
  $p = (p_0, p_1)$.
\item Let $C := \{$ready cells$\} \setminus \{f\}$.  If $C = \emptyset$
  the algorithm fails; otherwise, choose
  \[
    (r^*, c^*) := \arg\max_{(r,c) \in C}
       \Bigl( b_c, \; -r, \; -c \Bigr),
  \]
  with lexicographic comparison (priority: largest $b_c$, then smallest
  $r$, then smallest $c$).
\item Set $T(r^*, c^*) := j$, $b_{c^*} := b_{c^*} - 1$,
  $p := (r^*, c^*)$.
\end{enumerate}
\end{definition}

\begin{remark}\label{rem:why}
The selection rule prefers cells in the \emph{column with the most
remaining cells}.  Intuitively, this keeps all columns ``balanced'' so
that no single column gets ahead and forces a chain of consecutive
same-column placements (which would produce column descents).
\end{remark}

\section{Main result}\label{sec:main}

\begin{theorem}[Main]\label{thm:main}
For every partition $\lambda = (\hatl, 1^q)$ with $\hatl_i \ge 2$ for
all $i$, $\ell(\hatl) \ge 2$, $q \ge 1$, and $\tau(\hatl) = 0$, with
$|\lambda| \le 20$, Algorithm~N produces a standard Young tableau
$T^*$ of shape $\lambda$ satisfying conditions (i) and (ii) of
Conjecture~\ref{conj:main}.
\end{theorem}

\begin{proof}[Proof (computational verification)]
By exhaustive enumeration over all $\tau = 0$ shapes with
$|\lambda| \le 20$.  There are $1525$ such shapes.  For each, we
execute Algorithm~N and verify that
\begin{itemize}
\item the output $T$ is a standard Young tableau of shape $\lambda$
  (entries $1, \ldots, n$ each occurring once, rows and columns
  strictly increasing);
\item $T(0,0) = 1$ and $T(0,1) = 2$;
\item $T$ has no column descent.
\end{itemize}
The verification script
\verb|verify_construction_N.py| reports
\textbf{$1525/1525$ successful} (all pass).
\end{proof}

Combined with the previously-known brute-force verification for
$|\lambda| \le 14$ and the empirical exhaustive search, this gives
strong evidence for the conjecture.  We now sketch what would be
needed for a structural proof.

\section{Structural lemmas towards a full proof}\label{sec:structural}

\subsection{The SYT validity is automatic}

\begin{lemma}\label{lem:syt-valid}
For any partition $\lambda$ and any input, Algorithm~N either fails
(returns no tableau) or produces a valid SYT of shape $\lambda$.
\end{lemma}

\begin{proof}
At each step $j$, the algorithm places value $j$ in a ready cell.  A
ready cell $(r, c)$ has its left neighbor (if $c > 0$) and its top
neighbor (if exists) already filled, hence the row condition $T(r,c-1)
< T(r,c)$ and the column condition $T(r-1,c) < T(r,c)$ are
automatically satisfied.  Since each cell is placed exactly once and
values $1, \ldots, n$ are placed in order, the output is a valid SYT
whenever the algorithm completes.
\end{proof}

\subsection{Conditions (i) is automatic}

\begin{lemma}\label{lem:i-auto}
Whenever Algorithm~N completes successfully, the output $T$ satisfies
$T(0,0) = 1$ and $T(0,1) = 2$.
\end{lemma}

\begin{proof}
By construction (initialization step).
\end{proof}

\subsection{Condition (ii): the column descent question}

By construction, the only consecutive-value transitions $j, j+1$ that
need analysis are those produced by the greedy choice rule.  At step
$j+1$, the algorithm excludes the cell directly below the cell of $j$
from candidates.  Hence the only way a column descent at $j$ can occur
is if the forbidden cell was chosen anyway --- but the algorithm never
does this unless $C$ is empty (failure).  Therefore:

\begin{lemma}\label{lem:no-descent}
Whenever Algorithm~N completes successfully, the output $T$ has no
column descent.
\end{lemma}

\begin{proof}
By inspection of the algorithm: the cell directly below $p$ is in
$\{f\}$ and is excluded from $C$.  No cell from $C$ can equal $f$, so
the chosen cell is never directly below $p$.
\end{proof}

\subsection{The remaining gap: non-stuckness}

The only obstacle to a full proof is the following:

\begin{conjecture}\label{conj:not-stuck}
If $\lambda = (\hatl, 1^q)$ with $\hatl_i \ge 2$, $\ell(\hatl) \ge 2$,
$q \ge 1$, and $\tau(\hatl) = 0$, then Algorithm~N never fails (i.e.,
at every step $j$, the candidate set $C$ is non-empty).
\end{conjecture}

We have verified this conjecture exhaustively for $|\lambda| \le 20$
($1525$ shapes).  A structural proof would proceed by showing an
invariant about the column-fullness counters $b_c$ that prevents
stuckness.

\subsection{Counting necessary condition}\label{sec:counting}

The condition $\tau(\hatl) = 0$ is \emph{exactly} the necessary count:

\begin{lemma}\label{lem:tau0-necessary}
If a column-descent-free SYT $T^*$ of shape $\lambda = (\hatl, 1^q)$
exists, then $\tau(\hatl) = 0$, equivalently $q \le M + 1$.
\end{lemma}

\begin{proof}
Column $0$ of $\lambda$ has $a_0 = L + q$ cells in rows
$0, 1, \ldots, L+q-1$.  The labels in column $0$ form an increasing
sequence $v_0 < v_1 < \cdots < v_{L+q-1}$ with $v_i \in \{1, \ldots,
n\}$.  Column-descent-free means $v_{i+1} - v_i \ge 2$ for each $i$.
Hence
\[
  v_{L+q-1} \ge v_0 + 2(L+q-1).
\]
But $v_0 = 1$ and $v_{L+q-1} \le n = 2L + M + q$.  Thus
$2(L+q-1) \le 2L + M + q - 1$, i.e., $q \le M + 1$.
\end{proof}

So $\tau = 0$ is also the necessary condition for the conjecture to
hold; we are aiming to prove sufficiency.

\section{Status and discussion}\label{sec:status}

The conjecture is proved up to $|\lambda| \le 20$ via the explicit
deterministic Algorithm~N together with computational verification
(Theorem~\ref{thm:main}).  The full structural proof reduces to
Conjecture~\ref{conj:not-stuck} (non-stuckness of the greedy), which
we have verified extensively but not proved in closed form.

\subsection{What the algorithm does in special cases}

\begin{itemize}
\item For $\hatl$ with no length-$3$ row consecutive to another
  length-$\ge 3$ row, Algorithm~N reduces to the row-superstandard
  fill of $\hatl$ followed by tail, modulo small adjustments.
\item For multi-length-$3$-row cases (the gap in
  Proposition~4.8), Algorithm~N produces the same kind of clever
  interleaving as the manual examples constructed in the converse
  paper.
\item For tail-heavy cases ($q$ near $M+1$), the algorithm
  alternates tail cells with column-$\ge 3$ cells of $\hatl$,
  precisely the original ``interleaved $T^*$'' construction.
\end{itemize}

\subsection{Towards proving non-stuckness}

A key observation: the threshold $\tau = 0$ is equivalent to the
inequality
\[
  a_0 \le 1 + \sum_{c \ge 1} a_c
\]
(since $a_0 = L + q$, $\sum_{c \ge 1} a_c = L + M$, and $q \le M+1$).
We conjecture (but do not yet prove) that an invariant of the form
\begin{equation}\label{eq:invariant}
  b_0 \le 1 + \sum_{c \ge 1} b_c
\end{equation}
together with appropriate ``smaller column'' analogues is preserved
by Algorithm~N at every step, and that this invariant suffices to
prove $|C| \ge 1$ at every step.

\subsection{Verification artifacts}

\begin{itemize}
\item Algorithm implementation:
  \verb|/home/clio/projects/scratch/2026-05-16-existence-T-star/greedy10.py|,
  function \verb|construction_N|.
\item Verification script:
  \verb|/home/clio/projects/scratch/2026-05-16-existence-T-star/verify_construction_N.py|.
\item Brute-force confirmation:
  \verb|/home/clio/projects/scratch/2026-05-17-converse-prove/verify_existence_all.py|
  (independently confirms existence by backtracking search up to
  $|\lambda| \le 14$).
\end{itemize}

\end{document}
